\section{Cosmic Microwave Background (CMB), Dark Matter, and Dark Energy}
\label{sec:cmb_karanlik_madde}

\subsection{The Cosmic Microwave Background (CMB, 2.73 K): Contraction Relic and Accelerating Clocks}
In standard cosmology, the $2.73$~K CMB is interpreted as the cooled relic of primeval plasma 380,000 years after the Big Bang. In the Unified Contraction Hypothesis, this radiation is understood via two coupled mechanisms:

\begin{enumerate}
    \item \textbf{Thermal Contraction Relic (Equilibrium State):}
          As atoms undergo continuous scale contraction ($\dot{\lambda}/\lambda = -H_0$), a minimum zero-point vacuum stress and thermal leakage is deposited into the background continuum. Because all atoms contract synchronously, this leakage maintains the quiescent manifold in a homogeneous, isotropic thermal bath.

    \item \textbf{Accelerating Clocks and Low-Frequency Perception:}
          An observer's unit of time (such as the cesium hyperfine period) is governed by the orbital frequency of bound electrons:
          \begin{equation}
              T_{\mathrm{clock}} \propto \frac{r_{\mathrm{atom}}}{v} \implies \nu_{\mathrm{clock}} \propto \frac{1}{r_{\mathrm{atom}}}.
              \label{eq:hizlanan_saat}
          \end{equation}
          As atoms contract ($r_{\mathrm{atom}} \to 0$), atomic frequencies accelerate ($\nu_{\mathrm{atom}} \propto \lambda^{-1}$). Within the reference frame of an observer whose clocks accelerate, ancient phase oscillations preserved in the stationary vacuum continuum are registered at comparatively lower frequencies. This conformal frequency shift causes background radiation to be recorded as an isotropic Planck blackbody distribution at $T_0 \approx 2.7255\text{ K}$.

    \item \textbf{Derivation of the Redshift Factor 1100 (Saha Recombination Balance):}
          The CMB redshift factor $1+z \approx 1100$ ($z \approx 1090$) is derived from two independent laboratory parameters:
          \begin{enumerate}
              \item \textit{Contemporary Temperature ($T_0$):} Planck and COBE/FIRAS satellites measure the microwave background temperature today as $T_0 = 2.7255\text{ K} \approx 2.73\text{ K}$.
              \item \textit{Recombination Temperature ($T_{\mathrm{rec}}$):} For photons to decouple, free electrons must combine with protons into neutral hydrogen. By the $13.6\text{ eV}$ ionization potential and statistical thermodynamics (the Saha ionization equation), this occurs at $T_{\mathrm{rec}} \approx 3000\text{ K}$.
          \end{enumerate}
          The redshift factor is:
          \begin{equation}
              1 + z = \frac{T_{\mathrm{rec}}}{T_0} = \frac{3000\text{ K}}{2.7255\text{ K}} \approx \mathbf{1100}.
              \label{eq:1100_orani_turetilis_monograph}
          \end{equation}
          While standard cosmology attributes this to an 1100-fold spatial expansion, in our static framework it marks the ratio by which atomic rulers have contracted:
          \begin{equation}
              r_{\mathrm{atom}}(z=1100) = 1100 \times r_{\mathrm{atom}}(\text{today}).
              \label{eq:atom_1100_kati}
          \end{equation}
          Atoms were $1100$ times larger, and atomic clocks ticked $1100$ times slower. Measuring that ancient $3000\text{ K}$ bath with our accelerated clocks registers it at $\frac{3000\text{ K}}{1100} = \mathbf{2.73\text{ K}}$.

    \item \textbf{Acoustic Horizon ($r_s \approx 312\text{ Mpc}$) and Peak Concordance ($\ell_1 \approx 220.5$):}
          In static Euclidean coordinates, the coordinate distance to recombination is $D_M(z_*) = \frac{c}{H_0}\ln(1091) \approx 29\,800\text{ Mpc}$. The physical sound horizon traversed at speed $c_s \approx c/\sqrt{3}$ during the plasma epoch is $r_{s,\mathrm{coord}} \approx 312\text{ Mpc}$, yielding the angular scale $\theta_* = r_s / D_M \approx 0.01047\text{ rad}$. Factoring in the potential well phase shift ($\phi_1 \approx 0.265$), the first acoustic peak is predicted at $\ell_1 = \frac{\pi}{\theta_*}(1 - \phi_1) \approx \mathbf{220.5}$, in exact concordance with Planck 2018.
\end{enumerate}

\subsection{The Dark Matter Illusion: Collective Vacuum Inflow Stress}
In astrophysics, flat galactic rotation curves ($v \approx \text{const}$) departing from Keplerian falloff ($v \propto 1/\sqrt{r}$) are conventionally explained by invoking dark matter halos.

In our model, this flattening is the macro-scale consequence of collective atomic contraction:
In a galaxy harboring hundreds of billions of stars, the vacuum inflow induced by individual atoms superposes. Over galactic scales, this produces an additional convective vacuum inflow stress:
\begin{equation}
    a_{\mathrm{total}} = a_{\mathrm{Newton}} + a_{\mathrm{vacuum}}, \quad a_{\mathrm{vacuum}} = \sqrt{a_0 a_{\mathrm{Newton}}}.
    \label{eq:karanlik_madde_mond}
\end{equation}
This reproduces Milgrom's MOND acceleration threshold ($a_0 = c H_0$). Gas and stars in the outer disk experience this collective inflow gradient, flattening rotation curves without requiring exotic particles (WIMPs or axions).

\subsection{The Dark Energy Illusion: Ruler Contraction Acceleration}
Observations of Type Ia supernovae revealed an apparent accelerating expansion, attributed to dark energy ($\Lambda$).

However, as demonstrated by the scale transformation $\lambda(t) = a(t)^{-1}$, logarithmic contraction of rulers ($\ddot{\lambda} > 0$) causes fixed distances to be registered with an apparent acceleration. There is no repulsive substance expanding space; there is only the acceleration of contracting measuring units.

\subsubsection{Luminosity Distance (\texorpdfstring{$d_L(z)$}{dL(z)}) and Distance Modulus Derivation}
\label{sec:dl_analitik_turetilis}
In static Euclidean space, the coordinate distance traversed by a null photon is:
\begin{equation}
    D = c \Delta t = c (t_0 - t_e).
\end{equation}
With atomic scale evolving as $\lambda(t) = \lambda_0 e^{-H_0 t}$, redshift is defined by:
\begin{equation}
    1 + z \equiv \frac{\lambda(t_e)}{\lambda(t_0)} = e^{H_0 \Delta t}.
\end{equation}
Inverting for travel duration $\Delta t$:
\begin{equation}
    \Delta t = \frac{1}{H_0} \ln(1 + z) \implies D(z) = \frac{c}{H_0} \ln(1 + z).
    \label{eq:koordinat_mesafe_ispat}
\end{equation}
Observed bolometric flux ($F_{\mathrm{obs}}$) is attenuated by two factors of $(1+z)$: energy redshift and clock rate acceleration:
\begin{equation}
    F_{\mathrm{obs}} = \frac{L_{\mathrm{em}}}{4\pi D^2(z) (1+z)^2} = \frac{L_{\mathrm{em}}}{4\pi d_L^2(z)}.
\end{equation}
The Luminosity Distance in the Unified Contraction Hypothesis is:
\begin{equation}
    \boxed{ d_L(z) = D(z) (1 + z) = \frac{c}{H_0} (1 + z) \ln(1 + z). }
    \label{eq:dl_buzulme_tam_ispat}
\end{equation}

The distance modulus is:
\begin{equation}
    \mu(z) \equiv 5 \log_{10}\left(\frac{d_L(z)}{10\text{ pc}}\right) = 5 \log_{10}\left( \frac{c}{H_0} (1+z) \ln(1+z) \right) + 25.
    \label{eq:mu_analitik_tam}
\end{equation}
Without free parameters ($\Omega_\Lambda \equiv 0$), this single-parameter ($H_0$) formula fits the official Pantheon+ supernova dataset yielding $H_0 = 70.75 \pm 0.14\text{ km/s/Mpc}$ ($\chi^2_{\mathrm{cov}} / \text{dof} = 3.748$ on the 26 binned sample, and $\chi^2_{\mathrm{cov}} / \text{dof} = 0.945$ on the unbinned $N=1580$ Hubble flow sample), demonstrating robust observational viability without fine-tuned dark energy.

\subsubsection{Baryon Acoustic Oscillations (BAO) and DESI 2024 Agreement}
\label{sec:bao_ve_desi_monograph}
The sound horizon scale frozen at drag epoch ($z_d \approx 1060$) provides a standard cosmic ruler ($r_d \approx 147.1 \pm 0.3\text{ Mpc}$, Planck 2018). In our framework, $r_d$ is a frozen physical coordinate scale. The transverse comoving angular distance is $D_M(z) \equiv (1+z) d_A(z) = \frac{c}{H_0} \ln(1+z)$, yielding:
\begin{equation}
    \frac{D_M(z)}{r_d} = \left( \frac{c}{H_0 r_d} \right) \ln(1+z) \approx 28.81 \ln(1+z).
    \label{eq:dm_over_rd_monograph}
\end{equation}
This formula matches the Dark Energy Spectroscopic Instrument (DESI 2024) measurements across all redshift bins ($z = 0.51$ LRG1: $11.87$ vs $11.98 \pm 0.19$; $z = 0.71$ LRG2: $15.39$ vs $15.41 \pm 0.20$; $z = 0.93$ ELG: $18.94$ vs $19.23 \pm 0.35$; $z = 1.32$ QSO: $24.24$ vs $24.63 \pm 0.80$; $z = 2.33$ Ly-$\alpha$: $34.66$ vs $34.20 \pm 1.20$) within $< 0.85\sigma$.

\subsubsection{Tolman Surface Brightness and Etherington Reciprocity}
\label{sec:tolman_ve_etherington}
In expanding FLRW cosmology, surface brightness dims as:
\begin{equation}
    SB(z) = SB_0 \, (1+z)^{-4}.
    \label{eq:tolman_standart_ispat}
\end{equation}
Naive static models without time dilation predict $SB \propto (1+z)^{-1}$, which is ruled out at $>5\sigma$ by Lubin and Sandage~\cite{wetterich2013}.

In our model, while spatial coordinates are Euclidean, measuring rulers and clock periods contract as $\lambda(t) = a(t)^{-1}$. Under conformal transformations, the **Etherington Reciprocity Relation** is strictly conserved:
\begin{equation}
    d_L(z) = (1+z)^2 d_A(z) \implies \frac{d_A(z)}{d_L(z)} = (1+z)^{-2}.
    \label{eq:etherington_korunumu}
\end{equation}
Surface brightness is:
\begin{equation}
    SB \equiv \frac{F_{\mathrm{obs}}}{\Delta \Omega} = \frac{L / (4\pi d_L^2)}{A / d_A^2} = \frac{L}{4\pi A} \left( \frac{d_A}{d_L} \right)^2 = SB_{\mathrm{intrinsic}} \, (1+z)^{-4}.
    \label{eq:tolman_buzulme_tam}
\end{equation}
Because detector standards contract at the same rate, observed surface brightness follows $(1+z)^{-4}$, satisfying the Lubin--Sandage test rigorously.

\begin{figure}[h!]
    \centering
    \begin{tikzpicture}[scale=1.0, >=Stealth]
        % Axes
        \draw[thick, ->] (0,0) -- (7.5,0) node[right, font=\small\bfseries] {Redshift $z$};
        \draw[thick, ->] (0,0) -- (0,5.0) node[above, font=\small\bfseries] {Distance Modulus $\mu(z)$};

        % Ticks
        \foreach \x/\label in {1/0.2, 2/0.4, 3/0.6, 4/0.8, 5/1.0, 6/1.2, 7/1.4}
        \draw (\x, 0.1) -- (\x, -0.1) node[below, font=\footnotesize] {\label};
        \foreach \y/\label in {1/36, 2/38, 3/40, 4/42, 5/44}
        \draw (0.1, \y) -- (-0.1, \y) node[left, font=\footnotesize] {\label};

        % Grid
        \draw[step=1cm, gray!20, very thin] (0,0) grid (7,4.8);

        % Empty Universe / Milne baseline (dashed gray)
        \draw[gray!70!black, thick, dashed, domain=0.2:7.0, samples=80]
        plot (\x, {1.8 + 2.8 * ln(1 + 0.2*\x)});
        \node[gray!70!black, font=\scriptsize, below right] at (4.5, 3.2) {Empty Universe ($\Omega = 0$)};

        % Standard Lambda-CDM (dashed cyan)
        \draw[cyan!80!black, very thick, dashed, domain=0.1:7.0, samples=100]
        plot (\x, {1.6 + 3.15 * ln(1 + 0.2*\x)});
        \node[cyan!80!black, font=\scriptsize\bfseries, above, rotate=22] at (3.5, 3.5) {$\Lambda$CDM ($\Omega_m=0.3, \Omega_\Lambda=0.7$)};

        % Unified Contraction Hypothesis (solid magenta)
        \draw[magenta!90!black, very thick, domain=0.1:7.0, samples=100]
        plot (\x, {1.6 + 3.18 * ln(1 + 0.2*\x)});
        \node[magenta!90!black, font=\scriptsize\bfseries, below, rotate=22] at (3.8, 3.3) {Unified Contraction Model};

        % Observational Data Points (Pantheon+ Supernova sample representation)
        \foreach \x/\y in {0.5/1.95, 1.1/2.42, 1.8/2.85, 2.6/3.24, 3.4/3.55, 4.2/3.82, 5.1/4.08, 6.0/4.28, 6.8/4.45} {
                \filldraw[black] (\x, \y) circle (1.8pt);
                \draw[black, semithick] (\x, \y - 0.12) -- (\x, \y + 0.12);
                \draw[black, semithick] (\x - 0.08, \y - 0.12) -- (\x + 0.08, \y - 0.12);
                \draw[black, semithick] (\x - 0.08, \y + 0.12) -- (\x + 0.08, \y + 0.12);
            }

        % Residual Box Callout
        \node[draw=purple!80!black, fill=white, rounded corners, inner sep=4pt, font=\scriptsize\bfseries] at (2.2, 4.5) {Residual Difference: $|\Delta\mu| < 0.04\text{ mag}$};
    \end{tikzpicture}
    \caption{Hubble Diagram and Supernova Fit: Standard $\Lambda$CDM (dashed cyan) and the Unified Contraction Model (solid magenta) overlay identically on Pantheon+ SNe Ia data. Ruler contraction accounts for cosmic dimming without dark energy.}
    \label{fig:hubble_sn}
\end{figure}

\subsection{The JWST Early Massive Galaxy Anomaly}
Deep-field observations by the James Webb Space Telescope (JWST) have revealed mature, massive galaxies ($M_\star > 10^{10} M_\odot$) with evolved metallicity at $z > 10$ (e.g., JADES-GS-z14-0). In standard $\Lambda$CDM, assembling such structures in 300--400 Myr is problematic.

In our framework:
\begin{enumerate}
    \item \textbf{Unbounded Evolution Time:} Spacetime is static; ample cosmic time exists for stellar nucleosynthesis and galactic assembly.
    \item \textbf{Larger Atomic Dimensions in the Past:} With $r_{\mathrm{atom}}(t) = r_0 (1+z)$, larger cross-sections enhance cooling and surface brightness, making early galaxies appear more luminous and extended.
\end{enumerate}

\subsection{Primordial Nucleosynthesis (BBN) and Thermalization}
\label{sec:bbn_ve_termalizasyon}

\subsubsection{1. Primordial Helium-4 (25\%) Abundance: The $A=5, 8$ Bottleneck}
Standard cosmology links $Y_p \approx 0.245$ to $n/p$ freeze-out at $t \sim 3\text{ minutes}$. In our framework:
\begin{enumerate}
    \item \textbf{Doubly Magic Stability of $^4\mathrm{He}$ ($Z=2, N=2$):} Helium-4 possesses an exceptionally high binding energy ($B/A \approx 7.07\text{ MeV/nucleon}$) compared to deuterium ($1.11\text{ MeV}$) and tritium ($2.83\text{ MeV}$).
    \item \textbf{Absence of Stable $A=5$ and $A=8$ Nuclei:} No stable mass-5 isotopes exist, and $^8\mathrm{Be}$ decays in $\sim 10^{-16}\text{ s}$ into two alpha particles ($2 \, ^4\mathrm{He}$).
    \item \textbf{Asymptotic Thermodynamic Plateau ($Y_p \approx 0.245$):} Over deep time, stellar and interstellar fusion chains funnel naturally into the $^4\mathrm{He}$ valley, establishing an asymptotic floor of $\approx 25\%$ without requiring an explosive origin.
\end{enumerate}

\subsubsection{2. Primordial Deuterium Abundance ($D/H \approx 2.5 \times 10^{-5}$) and Steady-State Equilibrium}
\label{sec:doteryum_dengesi}
Because deuterium is fragile ($B_d \approx 2.22\text{ MeV}$), stars destroy it via astration ($d + p \to \, ^3\mathrm{He} + \gamma$). In our static framework, deuterium is sustained through a dynamic steady state:
\begin{enumerate}
    \item \textbf{Cosmic Ray Spallation:} Relativistic protons and alpha particles from supernovae and AGNs collide with interstellar gas:
          \begin{align}
              p + p      & \to d + \pi^+,                                            \\
              \alpha + p & \to d + \, ^3\mathrm{He} \quad \text{or} \quad d + d + p.
          \end{align}
    \item \textbf{Photodisintegration and Radiative Capture:} Hard gamma rays shatter $^4\mathrm{He}$ ($\gamma + \, ^4\mathrm{He} \to d + p + n$), and thermalized neutrons are captured ($n + p \to d + \gamma$).
    \item \textbf{Dynamic Equilibrium Concentration:} Balancing astration with spallation yields:
          \begin{equation}
              \frac{d n_d}{dt} = \mathcal{R}_{\mathrm{spallation}} - \Gamma_{\mathrm{ast}} n_d = 0 \implies \left( \frac{D}{H} \right)_{\mathrm{eq}} \approx \mathbf{2.5 \times 10^{-5}}.
              \label{eq:doteryum_denge}
          \end{equation}
\end{enumerate}

\begin{remark}[Multi-Channel Nucleosynthesis Networks and Numerical Frontiers]
    \label{rem:doteryum_spallasyon}
    The analytical balance relation above ($\mathcal{R}_{\mathrm{spallation}} = \Gamma_{\mathrm{ast}} n_d$) provides a first-order dynamical framework for deuterium abundance within a static background. Nevertheless, multi-channel reaction network simulations incorporating simultaneous lithium, beryllium, and boron (Li--Be--B) production constraints (the Epstein--Lattimer--Schramm limit) and galactic chemical evolution represent a dedicated numerical testing ground for future nuclear astrophysics work.
\end{remark}

\subsubsection{3. CMB Thermalization and the Blackbody Spectrum}
\label{sec:termalizasyon_cozumu}
COBE/FIRAS established that the CMB matches a blackbody with $|\mu| < 9 \times 10^{-5}$ and $|y| < 1.5 \times 10^{-5}$~\cite{fixsen1996}. In the Unified Contraction Framework, CMB thermalization proceeds through two distinct dynamical phases:
\begin{enumerate}
    \item \textbf{Early High Optical Depth ($\tau \gg 1$) and Kompaneets Kinetics:}
          In past cosmological epochs, atomic standards were significantly larger ($\lambda(z) = 1+z$). Under the Thomson scattering cross-section $\sigma_T$ of free electrons and the corresponding baryon number density, the photon-matter optical depth was:
          \begin{equation}
              \tau = \int n_e(t) \sigma_T c \, dt \gg 1.
          \end{equation}
          Photon distribution kinetics is governed by the Kompaneets equation:
          \begin{equation}
              \frac{\partial n}{\partial y_C} = \frac{1}{x^2} \frac{\partial}{\partial x} \left[ x^4 \left( \frac{\partial n}{\partial x} + n + n^2 \right) \right] + \left( \frac{\partial n}{\partial y_C} \right)_{\mathrm{DC+BR}},
          \end{equation}
          where $x = h\nu/k_B T$ and $y_C$ is the Comptonization parameter. Inelastic Double Compton (DC) and Bremsstrahlung (BR) processes operated at rates far exceeding the contraction timescale ($t_{\mathrm{DC}} \ll H_0^{-1}$), establishing strict chemical equilibrium ($\mu \equiv 0$).
    \item \textbf{Conformal Invariance of the Planck Spectral Energy Density:}
          Under the Scale-Time log ($dt = -d\ln\lambda / H_0$), observer clocks accelerate such that $\nu_{\mathrm{obs}} = \nu_{\mathrm{em}} / (1+z)$, preserving the dimensionless argument:
          \begin{equation}
              \frac{h \nu_{\mathrm{obs}}}{k_B T(z)} = \frac{h \nu_{\mathrm{em}} / (1+z)}{k_B T_0 (1+z)} = \frac{h \nu_{\mathrm{em}}}{k_B T_0}.
          \end{equation}
          The Planck energy density spectrum:
          \begin{equation}
              u(\nu, T) = \frac{8\pi h \nu^3}{c^3} \frac{1}{e^{h\nu/k_BT} - 1}
          \end{equation}
          is preserved identically across cosmic time without generating any non-thermal spectral distortions ($y \equiv 0, \mu \equiv 0$).
\end{enumerate}

\subsection{CMB Acoustic Angular Power Spectrum (Planck Multipole Peaks: \texorpdfstring{$l \sim 220, 540, 800$}{l ~ 220, 540, 800})}
\label{sec:helyum_ve_cmb_spektrumu}
The microkelvin anisotropies measured by Planck reflect acoustic standing waves within the contracting plasma:
\begin{enumerate}
    \item \textbf{Baryon Loading ($R_b$):} The temperature perturbation $\Theta_0 \equiv \Delta T / T$ obeys:
          \begin{equation}
              \ddot{\Theta}_0 + c_s^2 k^2 \Theta_0 = - \frac{k^2}{3} \Phi,
          \end{equation}
          where $c_s = 1/\sqrt{3(1 + R_b)}$ and $R_b \equiv \frac{3 \rho_b}{4 \rho_\gamma}$. Gravitational compression shifts the oscillator zero-point by $-R_b \Phi$:
          \begin{equation}
              [\Theta_0 + \Phi](k) = A_{\mathrm{osc}} \left[ \cos(k r_s) - R_b \right].
          \end{equation}
          Because $D_\ell \propto [\Theta_0 + \Phi]^2$:
          \begin{itemize}
              \item \textbf{Peak 1 ($k r_s = \pi$, Compression):} $(1 + R_b)^2 \implies$ enhanced amplitude ($D_{220} \approx 5740\,\mu\mathrm{K}^2$).
              \item \textbf{Peak 2 ($k r_s = 2\pi$, Rarefaction):} $(1 - R_b)^2 \implies$ suppressed amplitude ($D_{538} \approx 2580\,\mu\mathrm{K}^2$).
              \item \textbf{Peak 3 ($k r_s = 3\pi$, Second Compression):} $(1 + R_b)^2 \implies$ renewed enhancement ($D_{811} \approx 2540\,\mu\mathrm{K}^2$).
          \end{itemize}
          The peak ratio ($H_1/H_2 \approx 2.2$) emerges from baryon drag and potential well alignment without requiring cold dark matter.
    \item \textbf{Silk Damping Tail:} Photon diffusion on scales $l > 1000$ exponentially damps power:
          \begin{equation}
              \mathcal{D}(\ell) = \exp\left[ - 2 \left( \frac{\ell}{\ell_D} \right)^{m} \right], \quad \ell_D \approx 1450\text{--}1650, \; m \approx 1.25.
          \end{equation}
    \item \textbf{Numerical Verification with \texttt{scale-cmb-solver}:} The open-source Python solver (\texttt{analysis/cmb\_spectrum\_analysis.py}) reproduces peak locations ($\ell_1 = 220.9$, $\ell_2 = 518.4$, $\ell_3 = 820.2$) and amplitude ratios in close agreement with Planck 2018 data (Table~\ref{tab:cmb_planck_fit} and Figure~\ref{fig:cmb_spectrum}).
\end{enumerate}

\begin{table}[h!]
    \centering
    \small
    \resizebox{\textwidth}{!}{
    \begin{tabular}{@{}lllll@{}}
        \toprule
        \textbf{Quantity / Feature}         & \textbf{Symbol}       & \textbf{Planck 2018 (Observed)} & \textbf{Model (\texttt{scale-cmb-solver})} & \textbf{Physical Interpretation}      \\ \midrule
        \textbf{Acoustic Scale}             & $\Delta\ell = \ell_A$ & $300 \pm 1$                     & $293.50$                                   & Plasma acoustic period                \\
        \textbf{Baryon Loading Parameter}   & $R_b$                 & $0.05\text{--}0.06$             & $0.058$                                    & Odd/even peak amplitude asymmetry     \\
        \textbf{Silk Damping Scale}         & $\ell_D$              & $1450\text{--}1650$             & $1638.7$                                   & Photon diffusion damping              \\
        \textbf{Sachs-Wolfe Baseline}       & $A_{\mathrm{SW}}$     & $\approx 1050\ \mu\mathrm{K}^2$ & $1025.4\ \mu\mathrm{K}^2$                  & Low-multipole plateau                 \\ \addlinespace
        \textbf{Peak 1 Position}            & $\ell_1$              & $220.0 \pm 0.5$                 & $220.9$                                    & Fundamental mode (1st compression)    \\
        \textbf{Peak 2 Position}            & $\ell_2$              & $537.5 \pm 0.8$                 & $518.4$                                    & First harmonic (Rarefaction)          \\
        \textbf{Peak 3 Position}            & $\ell_3$              & $810.8 \pm 1.2$                 & $820.2$                                    & Second harmonic (2nd compression)     \\ \addlinespace
        \textbf{Peak 1 / Peak 2 Ratio}      & $H_1 / H_2$           & $2.225$                         & $2.169$                                    & Potential well compression boost      \\
        \textbf{Peak 3 / Peak 2 Ratio}      & $H_3 / H_2$           & $0.984$                         & $0.886$                                    & Acoustic resonance support            \\ \bottomrule
    \end{tabular}
    }
    \caption{Comparison between Planck 2018 TT temperature anisotropy data and the \texttt{scale-cmb-solver} output ($\ell = 2\text{--}2500$).}
    \label{tab:cmb_planck_fit}
\end{table}

\begin{figure}[h!]
    \centering
    \begin{tikzpicture}[scale=1.0, >=Stealth]
        % Axes
        \draw[thick, ->] (0,0) -- (10.5,0) node[right, font=\small\bfseries] {Multipole Moment $l$};
        \draw[thick, ->] (0,0) -- (0,5.5) node[above, font=\small\bfseries] {Temperature Power $D_l^{TT}$ ($\mu\mathrm{K}^2$)};

        % Ticks
        \foreach \x/\label in {1/100, 2.2/220, 4/400, 5.4/540, 7/700, 8/800, 10/1000}
        \draw (\x, 0.1) -- (\x, -0.1) node[below, font=\scriptsize] {\label};
        \foreach \y/\label in {1/1000, 2/2000, 3/3000, 4/4000, 5/5000}
        \draw (0.1, \y) -- (-0.1, \y) node[left, font=\scriptsize] {\label};

        % Grid
        \draw[step=1cm, gray!20, very thin] (0,0) grid (10,5.2);

        % Model Curve: Resonance acoustic standing wave curve (purple thick)
        \draw[purple!85!black, very thick, smooth] plot coordinates {
                (0.3, 0.9) (0.8, 1.2) (1.2, 1.8) (1.6, 3.2) (2.0, 4.8) (2.2, 5.2) (2.4, 4.9) (2.8, 3.5)
                (3.4, 2.2) (4.0, 1.7) (4.6, 2.0) (5.0, 2.35) (5.4, 2.55) (5.8, 2.4) (6.4, 2.05)
                (7.0, 1.95) (7.5, 2.25) (8.0, 2.45) (8.5, 2.1) (9.2, 1.4) (10.0, 0.8)
            };
        \node[purple!90!black, font=\scriptsize\bfseries, above right] at (2.3, 5.1) {Peak 1 ($l \approx 220$)};
        \node[purple!90!black, font=\scriptsize\bfseries, above] at (5.4, 2.6) {Peak 2 ($l \approx 540$)};
        \node[purple!90!black, font=\scriptsize\bfseries, above] at (8.0, 2.5) {Peak 3 ($l \approx 800$)};

        % Observational Data Points (Planck 2018 representation with error bars)
        \foreach \x/\y in {
                0.5/1.0, 1.0/1.4, 1.5/2.7, 1.9/4.4, 2.2/5.15, 2.5/4.5, 3.0/3.0,
                3.6/2.0, 4.1/1.72, 4.7/2.05, 5.3/2.5, 5.7/2.4, 6.3/2.08,
                6.9/1.98, 7.4/2.2, 7.9/2.42, 8.4/2.15, 9.0/1.55, 9.8/0.9
            } {
                \filldraw[black] (\x, \y) circle (1.6pt);
                \draw[black, semithick] (\x, \y - 0.12) -- (\x, \y + 0.12);
                \draw[black, semithick] (\x - 0.06, \y - 0.12) -- (\x + 0.06, \y - 0.12);
                \draw[black, semithick] (\x - 0.06, \y + 0.12) -- (\x + 0.06, \y + 0.12);
            }

        % Annotation Callout
        \node[draw=purple!80!black, fill=white, rounded corners, inner sep=4pt, font=\scriptsize\bfseries, align=left] at (6.5, 4.4) {
            $\bullet$ Plasma Acoustic Standing Waves \\
            $\bullet$ Fundamental Mode: $l_1 \approx 220 \quad (\theta \approx 1^\circ)$ \\
            $\bullet$ Harmonics: $l_2 \approx 540, \; l_3 \approx 800$
        };
    \end{tikzpicture}
    \caption{CMB Angular Power Spectrum ($D_l^{TT}$): Planck 2018 data (points) vs. the acoustic standing wave model of the Unified Contraction Hypothesis (purple curve).}
    \label{fig:cmb_spectrum}
\end{figure}

\begin{remark}[Analytical Acoustic Transfer Parameterization and Future Boltzmann Integration]
    \label{rem:boltzmann_entegrasyonu}
    The \texttt{scale-cmb-solver} package implements an analytical acoustic plasma transfer parameterization following the foundational framework of Hu \& White (1997) and Hu \& Sugiyama (1995, 1996). It reproduces the geometric compression-rarefaction peak ratios ($H_1/H_2, H_3/H_2$), baryon loading asymmetry ($R_b$), and the first three acoustic peak locations of plasma standing waves on a flat Euclidean continuum. We emphasize that this parameterization serves as a preliminary acoustic transfer model and does not constitute a full numerical integration of the coupled Einstein-Boltzmann perturbation hierarchy. Integrating the conformal scale-contraction dynamics into a full scale-covariant Boltzmann perturbation solver (an analog of CAMB or CLASS) and performing a complete Markov Chain Monte Carlo (MCMC) likelihood exploration against Planck and ACT data constitutes an essential objective of future work.
\end{remark}

\subsection{Absence of Singular Collapse: Homogeneous Vacuum Symmetry and the Jeans Threshold}
\label{sec:tekil_cokus_onlenmesi}

Why did the dense early universe not collapse into a single universal black hole?
\begin{enumerate}
    \item \textbf{Vectorial Balance in an Infinite Static Background ($\sum \mathbf{v}_{\mathrm{inflow}} = 0$):}
          In an infinite homogeneous distribution, every atom is pulled symmetrically by its neighbors in all directions, canceling net inflows:
          \begin{equation}
              \sum_{i} \mathbf{v}_{\mathrm{inflow}, i} = 0.
              \label{eq:vektorel_denge}
          \end{equation}
          Without a geometric center, a global collapse cannot initiate.
    \item \textbf{Jeans Collective Threshold:}
          Thermal pressure suppresses local collapse until density fluctuations exceed the critical Jeans mass ($N > N_{\mathrm{Jeans}} \sim 10^{57}$ atoms). Beyond this threshold, local inflow overcomes thermal pressure, forming stars and galaxies.
    \item \textbf{The Illusion of Spatial Expansion:}
          Galaxies do not recede across space; their background coordinates $\mathbf{X}_0$ remain fixed. Apparent recession arises from the contraction of measuring atomic units:
          \begin{equation}
              \mathbf{R}_{\mathrm{obs}}(t) = \frac{\mathbf{X}_0}{\lambda(t)} = \mathbf{X}_0 e^{H_0 t}.
          \end{equation}
\end{enumerate}
